\chapter{Introduction}
\section{Drug discovery and affinity prediction}
Developing a medicine involves identifying a biological target, finding molecules that affect it, and testing whether the effect is useful and safe. Laboratory screening is essential but can be costly and time-consuming. Computational methods can prioritize candidates and help researchers analyze large chemical and biological spaces; they complement rather than replace experiments.

Drug--target affinity describes how strongly a molecule binds a target protein. In this project, each example is a compound--protein pair and the output is a real-valued affinity target. A predicted value is a model estimate, not a measured binding result. DTA prediction can help prioritize experiments, while its usefulness depends on data quality, evaluation design, and generalization to new compounds or targets.

Traditional pipelines often use hand-crafted molecular descriptors and sequence descriptors, or treat the molecule and protein as fixed strings. Such representations can obscure molecular connectivity or local sequence patterns. A molecular graph naturally treats atoms as nodes and chemical bonds as edges. A protein sequence retains residue order and allows convolutional filters to detect local motifs. Neither representation alone provides a full physical model of a protein--ligand complex.

\section{Project objectives}
The project objectives were to implement the main MGraphDTA forward architecture for the Davis data, construct molecular graph and protein sequence inputs, train and evaluate a continuous affinity regressor, isolate the provided test fold from fitting and model selection, implement the project's evaluation metrics and Grad-AAM, and record additional protein-graph and structural investigations without conflating them with the primary model's results.

\chapter{Original Paper}
\section{Paper information and problem}
The original work is Ziduo Yang, Weihe Zhong, Lu Zhao, and Calvin Yu-Chian Chen, ``MGraphDTA: deep multiscale graph neural network for explainable drug--target binding affinity prediction,'' published in \emph{Chemical Science}, volume 13, pages 816--833 (2022), DOI 10.1039/D1SC05180F \cite{yang2022mgraphdta}. It addresses prediction of continuous drug--target affinity and the limitation of shallow graph encoders in representing molecular structure at multiple depths.

\section{Original architecture}
The paper's central idea is to encode a compound as a molecular graph with a deep, densely connected multiscale graph neural network (MGNN), encode a target amino-acid sequence with a multiscale convolutional neural network (MCNN), combine the learned representations, and regress affinity with an MLP. It describes 27 graph-convolutional architectural units in MGNN, three protein convolution branches at different depths, and Grad-AAM for atom-level visual interpretation. Sequence characters are mapped to integer tokens; molecular atoms and bond connectivity form the graph. The authors report benchmark experiments over seven DTA datasets, including Davis, and compare using MSE, CI, and $r_m^2$.

The original paper reports Davis MSE $0.207\pm0.001$, CI $0.900\pm0.004$, and $r_m^2=0.710\pm0.005$. Those are paper results under its benchmark protocol, not values generated by this repository. The paper's seven datasets and tasks are summarized in Table~\ref{tab:paper-datasets} \cite{yang2022mgraphdta}. Its regression setup reports Adam with learning rate 0.0005 and batch size 512; hyperparameters were optimized on Davis with five-fold cross-validation using Optuna, then kept fixed for other datasets. The paper also examines random, orphan-drug, and cluster-based splits and an independent test procedure for a filtered Davis evaluation. These choices differ from this project's single-seed split-and-validation workflow. The full numerical results on all seven datasets are not reproduced here.

\begin{table}[H]\centering\small\caption{Datasets used in the original paper}\label{tab:paper-datasets}\begin{tabular}{llrrr}\toprule Dataset & Task & Compounds & Proteins & Interactions\\\midrule Davis & Regression & 68 & 442 & 30,056\\Filtered Davis & Regression & 68 & 379 & 9,125\\KIBA & Regression & 2,111 & 229 & 118,254\\Metz & Regression & 1,423 & 170 & 35,259\\Human & Classification & 2,726 & 2,001 & 6,728\\C. elegans & Classification & 1,767 & 1,876 & 7,786\\ToxCast & Regression & 3,098 & 37 & 114,626\\\bottomrule\end{tabular}\end{table}

\section{Training, data, metrics, and explanation in the paper}
The paper defines the architecture and evaluates it on benchmark DTA datasets with repeated experimental reporting. The project's local paper copy is \texttt{bio\_paper.pdf}; bibliographic metadata and reported values are cross-checked against the publisher's article record. The paper uses distinct split settings across its evaluations, including five-fold random split experiments and orphan-drug/cluster-based split analyses; it is therefore inaccurate to reduce all its reported numbers to one split. This project does not reproduce those protocols. Grad-AAM weights activation channels by gradients of the scalar model output and maps a weighted activation sum to molecular atoms. It provides model attribution, not experimental interaction evidence.

\section{Paper, authors' code, and this implementation}
The paper is the source for the proposed MGNN/MCNN method, benchmark protocol descriptions, and paper-reported scores. The authors' public repository provides executable implementations and dataset-specific training/testing workflows; its regression entry points document a batch-size default of 512 \cite{mgraphdtaauthorscode}. This project implements its own data objects, model modules, training loop, validation selection, and evaluation-only checkpoint path in the files documented later in this report. It uses seed 42, divides the provided 25,046-row training fold into 22,541 optimization and 2,505 validation rows, and evaluates the separate 5,010-row test fold after restoring the best validation checkpoint. The project's test isolation and split are implementation choices, not a claim that the paper or the authors' code used this identical procedure.

\section{Original method components in more detail}
The original paper's drug input is a molecular graph whose node and adjacency representation supports message passing. MGNN is made deep through dense connections and multiscale graph-convolution blocks; the paper's 27-unit statement is its architectural layer convention. The protein input is a tokenized amino-acid sequence processed through convolution branches at increasing depths to capture different sequence-context scales. The two learned representations are combined and passed through an MLP to predict affinity. Grad-AAM uses output gradients to weight activation channels and produce per-atom visual attributions. The original work evaluates seven benchmark datasets and reports MSE, CI, and $r_m^2$ over repeated experiments. Its reported Davis values appear later in Table~\ref{tab:paper-vs-project}. Numeric dimensions from local source code are presented only in the implementation chapter; they are not presumed to be paper specifications when not explicitly verified.

\chapter{Dataset and Representations}
\section{Davis dataset}
The locally available raw Davis source files are \texttt{ligands\_can.txt}, \texttt{proteins.txt}, \texttt{Y}, and fold index files under \texttt{data/raw/davis/}. The verified dataset contains 30,056 interactions, 68 compounds, 442 target IDs, and 379 unique protein sequences. The counts are consistent with the project's dataset documentation and locally processed CSVs. Distinct target IDs can share a sequence. SMILES encode chemical connectivity; the protein input is the amino-acid sequence. The recorded affinity source is $K_d$ in nM.

The project transforms the Davis affinity values as
\begin{equation}\label{eq:pkd} y = 9-\log_{10}(K_d\text{ in nM}). \end{equation}
A smaller $K_d$ means less compound is needed to reach the measured binding equilibrium and thus stronger binding. Since the transform is decreasing in $K_d$, a higher transformed affinity corresponds to a smaller $K_d$. The values are regression targets, not model-generated labels.

\begin{table}[H]\centering\caption{Verified Davis dataset and split counts}\label{tab:data}\begin{tabular}{lr}\toprule Quantity & Count\\\midrule Interactions, total & 30,056\\Unique compounds & 68\\Target IDs & 442\\Unique protein sequences & 379\\Training fold used by project & 25,046\\Independent test fold & 5,010\\Optimization subset (90\% of training fold) & 22,541\\Validation subset (10\% of training fold) & 2,505\\\bottomrule\end{tabular}\end{table}

The project uses the provided train and test folds. Only the training fold is randomly split into optimization and validation subsets (seed 42). Validation selects the checkpoint and drives stopping; the independent test fold is evaluated after the best checkpoint is restored. This test isolation and validation procedure are project choices and are not identical to the paper's exact benchmark protocol.

\chapter{Implemented System Architecture}
\section{End-to-end architecture}
Figure~\ref{fig:overall} depicts the main implemented pathway. The system is a two-input regressor: drug features and target sequence features are encoded separately, concatenated, and passed to the prediction head.
\begin{figure}[H]\centering\begin{tikzpicture}[node distance=7mm and 13mm, box/.style={draw=navy,rounded corners,fill=lightblue,align=center,text width=35mm,minimum height=9mm},arrow/.style={-{Latex[length=2mm]},thick,draw=navy}]
\node[box] (s) {Drug SMILES}; \node[box,below=of s] (rd) {RDKit parsing}; \node[box,below=of rd] (g) {Molecular graph\\22-D atom features + edges}; \node[box,below=of g] (mg) {MGNN}; \node[box,below=of mg] (de) {96-D drug embedding};
\node[box,right=55mm of s] (p) {Protein amino-acid sequence}; \node[box,below=of p] (tok) {Tokenize, pad/truncate\\to 1,200}; \node[box,below=of tok] (mc) {128-D embedding + MCNN}; \node[box,below=of mc] (pe) {96-D protein embedding};
\node[box,below=20mm of de,xshift=27mm] (cat) {Concatenate: 192-D}; \node[box,below=of cat] (mlp) {MLP\\192 $\to$ 1024 $\to$ 1024 $\to$ 256 $\to$ 1}; \node[box,below=of mlp] (out) {Predicted affinity};
\draw[arrow] (s)--(rd);\draw[arrow](rd)--(g);\draw[arrow](g)--(mg);\draw[arrow](mg)--(de);\draw[arrow](p)--(tok);\draw[arrow](tok)--(mc);\draw[arrow](mc)--(pe);\draw[arrow](de)--(cat);\draw[arrow](pe.south)--++(0,-6mm)-|(cat.east);\draw[arrow](cat)--(mlp);\draw[arrow](mlp)--(out);\end{tikzpicture}\caption{Implemented MGraphDTA dataflow. The diagram represents the repository's current sequence-based protein branch.}\label{fig:overall}\end{figure}

\section{Molecular graph and atom features}
RDKit parses each SMILES. Atoms become graph nodes; each chemical bond becomes two directed edge-index entries (one in each direction). The model receives node features and connectivity. No separate bond-feature tensor is returned by current preprocessing, so the report does not attribute learned bond types to the input.

\begin{longtable}{p{0.08\linewidth}p{0.65\linewidth}p{0.12\linewidth}}\caption{All implemented atom feature dimensions}\label{tab:atomfeatures}\\\toprule Index & Feature & Dimensions\\\midrule\endfirsthead\toprule Index & Feature & Dimensions\\\midrule\endhead
1 & Atom symbol one-hot: H, C, N, O, F, Cl, S, Br, I & 9\\
2 & Atomic number & 1\\3 & Hydrogen-bond acceptor flag from RDKit chemical features & 1\\4 & Hydrogen-bond donor flag from RDKit chemical features & 1\\5 & Aromaticity flag & 1\\6 & Hybridization one-hot: SP, SP2, SP3 & 3\\7 & Total hydrogen count & 1\\8 & Explicit valence & 1\\9 & Formal charge & 1\\10 & Implicit valence & 1\\11 & Explicit hydrogen count & 1\\12 & Radical electron count & 1\\\midrule Total & & 22\\\bottomrule\end{longtable}

The one-hot groups contribute 12 dimensions (9 element and 3 hybridization); the remaining ten scalar properties bring the total to 22. Scalar values are emitted as floating-point node features. Atom-feature normalization beyond the code's conversion to floats is not applied by the described extraction function.

\section{MGNN architecture}
The initial graph convolution maps 22 atom channels to 32. Each of three dense blocks contains eight DenseLayer units. A DenseLayer applies a bottleneck GraphConv to an intermediate width of $2\times32=64$, then a second GraphConv emits 32 new channels; the output is concatenated with the layer input. Each transition maps the block output to half as many channels. Global mean pooling reduces a variable-size molecular graph to one vector, then a linear layer maps 228 channels to a 96-dimensional drug embedding.

\begin{table}[H]\centering\small\caption{MGNN channel progression}\label{tab:mgnn}\begin{tabular}{lrrl}\toprule Stage & Input channels & Output channels & Operation\\\midrule Initial conv0 & 22 & 32 & GraphConv + normalization/activation as implemented\\Dense block 1 (8 units) & 32 & 288 & Dense concatenation, growth 32\\Transition 1 & 288 & 144 & GraphConv, channel compression\\Dense block 2 (8 units) & 144 & 400 & Dense concatenation, growth 32\\Transition 2 & 400 & 200 & GraphConv, channel compression\\Dense block 3 (8 units) & 200 & 456 & Dense concatenation, growth 32\\Transition 3 & 456 & 228 & GraphConv, channel compression\\Readout & 228 & 96 & Global mean pooling + linear projection\\\bottomrule\end{tabular}\end{table}

The code documents three distinct depth counts. The paper-style count is 27: 24 DenseLayer units plus three transitions, excluding conv0. Counting conv0 gives 28 top-level architectural stages. Counting literal PyG GraphConv modules gives 52 because each DenseLayer has two GraphConv operations and conv0 and transitions each contribute one ($1+24\times2+3$). These counts refer to different objects and should not be conflated. The dense blocks preserve and concatenate earlier features, enabling information from multiple graph depths to reach later stages. Figure~\ref{fig:mgnn} summarizes the stage progression.
\begin{figure}[H]\centering\begin{tikzpicture}[node distance=5mm,box/.style={draw=navy,rounded corners,fill=lightblue,align=center,text width=95mm,minimum height=7mm},arr/.style={-{Latex},thick}]
\node[box](a){22 atom channels $\to$ conv0 $\to$ 32};\node[box,below=of a](b){Dense block 1: 8 units, growth 32 $\to$ 288};\node[box,below=of b](c){Transition 1: 288 $\to$ 144};\node[box,below=of c](d){Dense block 2: 8 units $\to$ 400};\node[box,below=of d](e){Transition 2: 400 $\to$ 200};\node[box,below=of e](f){Dense block 3: 8 units $\to$ 456};\node[box,below=of f](g){Transition 3: 456 $\to$ 228};\node[box,below=of g](h){Global mean pooling + FC $228\to96$};\foreach \x/\y in {a/b,b/c,c/d,d/e,e/f,f/g,g/h}{\draw[arr](\x)--(\y);}\end{tikzpicture}\caption{MGNN channel stages from the implementation. DenseLayer internals use bottleneck width 64 then add 32 channels.}\label{fig:mgnn}\end{figure}

\section{MCNN protein encoder}
Protein symbols use the project's 25-character vocabulary with index zero reserved for padding. Inputs are normalized to uppercase, mapped to integer IDs, and padded or truncated to 1,200 tokens. The trainable embedding width is 128. Three parallel branches contain one, two, or three Conv1D layers. Each layer uses 96 output channels, kernel size 3, stride 1, no padding, and ReLU. Adaptive max pooling returns 96 values per branch; concatenation gives 288 values and a final linear layer yields 96 dimensions. With a length-1,200 input the output lengths before pooling are 1,198, 1,196, and 1,194.
\begin{table}[H]\centering\caption{MCNN branches and shapes}\label{tab:mcnn}\begin{tabular}{lrrrr}\toprule Branch & Conv1D layers & Kernel/stride/padding & Channels & Length before pool\\\midrule 1 & 1 & 3 / 1 / 0 & 96 & 1198\\2 & 2 & 3 / 1 / 0 & 96 & 1196\\3 & 3 & 3 / 1 / 0 & 96 & 1194\\\bottomrule\end{tabular}\end{table}
The non-padding behavior follows the implementation. It should not be mistaken for an explicit numeric padding prescription in the paper text. Figure~\ref{fig:mcnn} shows the three branches.
\begin{figure}[H]\centering\begin{tikzpicture}[node distance=8mm and 9mm,box/.style={draw=navy,rounded corners,fill=lightblue,align=center,text width=31mm,minimum height=8mm},arr/.style={-{Latex},thick}]
\node[box](in){Tokens $[B,1200]$};\node[box,below=of in](emb){Embedding $128$};\node[box,below left=of emb](b1){Branch 1\\1 Conv1D};\node[box,below=of emb](b2){Branch 2\\2 Conv1D};\node[box,below right=of emb](b3){Branch 3\\3 Conv1D};\node[box,below=of b2](pool){Adaptive max pool\\each branch $96$};\node[box,below=of pool](cat){Concatenate $288$};\node[box,below=of cat](proj){Linear projection $288\to96$};\draw[arr](in)--(emb);\draw[arr](emb)--(b1);\draw[arr](emb)--(b2);\draw[arr](emb)--(b3);\draw[arr](b1.south)|-(pool.west);\draw[arr](b2)--(pool);\draw[arr](b3.south)|-(pool.east);\draw[arr](pool)--(cat);\draw[arr](cat)--(proj);\end{tikzpicture}\caption{MCNN branch structure and fusion. Each convolution uses kernel 3 and 96 channels.}\label{fig:mcnn}\end{figure}

\section{MLP regression head}
The 96-dimensional protein and 96-dimensional drug embeddings are concatenated into a 192-dimensional vector. The implemented head is $192\to1024\to1024\to256\to1$. Each of the three hidden linear layers is followed by ReLU and dropout with probability 0.1. The final scalar is an unconstrained regression output on the transformed affinity scale; it is not a class probability.
